% color and xcolor already loaded by eccv.sty
\usepackage{soul}
\usepackage{multirow}
% \usepackage{algorithm}    % uncomment when needed
% \usepackage{algorithmic}  % uncomment when needed
\usepackage{wrapfig}
\definecolor{darkred}{rgb}{0.7,0.1,0.1}
\definecolor{darkgreen}{rgb}{0.1,0.7,0.1}
\definecolor{cyan}{rgb}{0.7,0.0,0.7}
\definecolor{dblue}{rgb}{0.2,0.2,0.8}
\definecolor{maroon}{rgb}{0.76,.13,.28}
\definecolor{burntorange}{rgb}{0.81,.33,0}
\definecolor{tealblue}{rgb}{0.212,0.459, 0.533}
\definecolor{myyellow}{rgb}{0.8627451 , 0.75294118, 0.20784314}

\definecolor{mypink}{rgb}{0.93359375, 0.62109375, 0.83984375}

\definecolor{pp}{rgb}{0.43921569, 0.18823529, 0.62745098}
\definecolor{rr}{rgb}{0.5254902 , 0.00784314, 0.12941176}
\definecolor{bb}{rgb}{0.09019608, 0.23529412, 0.37647059}
\definecolor{yy}{rgb}{0.49803922, 0.3372549 , 0.0}
\definecolor{gg}{rgb}{0.02352941, 0.3372549 , 0.17647059}

%%% AF: see line at top of main file to show/hide notes
\ifdefined\ShowNotes
  \newcommand{\colornote}[3]{{\color{#1}\bf{#2: #3}\normalfont}}
\else
  \newcommand{\colornote}[3]{}
\fi

%%%% Add your commend here %%%%
\newcommand {\ray}[1]{\colornote{tealblue}{Ray}{#1}}
\newcommand {\tim}[1]{\colornote{burntorange}{Tim}{#1}}

\definecolor{lightred}{rgb}{0.9,0.4,0.4}
\newcommand{\lightred}[1]{{\color{lightred}{#1}}}
\newcommand{\darkgreen}[1]{{\color{darkgreen}{#1}}}
\newcommand{\lightblue}[1]{{\color{blue}{#1}}}

\definecolor{almostblack}{HTML}{111111}

\definecolor{darkpink}{rgb}{0.98,0.81,0.89}

\newcommand{\eat}[1]{} % ignore argument
\newcommand{\bsla}{\textit{Only$\gR$}\xspace}
\newcommand{\bslb}{\textit{No$F$}\xspace}

\definecolor{OursColor}{rgb}{0.83,0.89,0.94}
\newcommand{\ours}{\cellcolor{OursColor}}
\newcommand{\best}[1]{\textbf{#1}}
\newcommand{\second}[1]{\underline{#1}}


\definecolor{wolto}{rgb}{.75,0.75,0.75}

\newcommand{\centercell}[1]{\multicolumn{1}{c}{#1}}


\newcommand{\hlc}[2][yellow]{{%
    \colorlet{foo}{#1}%
    \sethlcolor{foo}\hl{#2}}%
}


%%%%%%%%%%%%%%%%%%%%%%%%

\newlength\savewidth\newcommand\shline{\noalign{\global\savewidth\arrayrulewidth\global\arrayrulewidth 1pt}\hline\noalign{\global\arrayrulewidth\savewidth}}

\definecolor{turquoise}{cmyk}{0.65,0,0.1,0.1}
\definecolor{purple}{rgb}{0.65,0,0.65}
\definecolor{darkgreen}{rgb}{0.0, 0.5, 0.0}
\definecolor{darkred}{rgb}{0.5, 0.0, 0.0}
\definecolor{darkblue}{rgb}{0.0, 0.0, 0.5}
\definecolor{blue}{rgb}{0.0, 0.0, 1.0}
\definecolor{orange}{rgb}{1.0,0.5,0.0}

%use this in running text
\newcommand{\on}{{\textit{on} }}
\newcommand{\off}{{\textit{off} }}

%use this in image captions
\newcommand{\ON}{{\textup{on} }}
\newcommand{\OFF}{{\textup{off} }}
\newcommand{\OFFCapital}{{\textup{Off} }}

\newcommand{\warn}{\color{red} \textbf{WARNING!} }
\newcommand{\warntwo}[1]{\color{red} \textbf{WARNING! #1}}
\newcommand{\warnoff}{\color{black}}


\newcommand{\todo}[1]{{\textbf{\color{blue}[todo: #1]}}}
\newcommand{\changed}[1]{{\color{blue}#1}}
\newcommand{\ch}[1]{{\color{orange}#1}}
\newcommand{\erase}[1]{\st{#1}}


\newcommand{\mbf}[1]{\mathbf{#1}}
\renewcommand{\vec}[1]{\mathbf{#1}}
\newcommand{\abs}[1]{\lvert #1 \rvert}

\definecolor{code_bg}{rgb}{0.95, 0.95, 0.95}

\newcommand{\nv}{NVidia\texttrademark\,}

\newif\ifproofread

\newcommand{\newcontent}[1]{%
\ifproofread
\textcolor{blue}{#1}%
\else
#1%
\fi
}

% \proofreadtrue
